\chapter{Which Film Is the Film}\label{ch:which-film}

The hero of the thriller at the origin of this book is named Theseus. The name recalls an old puzzle. The ship in which Theseus sailed was preserved by the Athenians, who replaced its planks one by one as they rotted, until none of the original timber remained, and philosophers asked whether it was still the same ship \cite{plutarch}. Selective cinema poses a stranger version of the question. It is not that one ship's planks are replaced over time. It is that several ships sail at once, sharing some planks and not others, and all of them carry the same name. Which is the ship?

\section{The question}

The question is not idle. Critics must know what they are reviewing. Archives must know what to preserve. Courts, in disputes about rights, defamation, or contracts, must know what a work contains. Audiences who argue about a film must know whether they are arguing about the same thing. In ordinary cinema the answer is conventional: the work is the released cut, and variants are versions of it. In selective cinema, where several realizations are released together, shown together, and none is necessarily primary, the convention does not apply.

\section{Four answers}

There are four candidate answers.

The first privileges one realization as the \emph{master} and treats the others as variants of it. This is the answer of the primary family in \Cref{ch:family}, and for some families it is correct: a work made as a feature with an explanatory version for students has a clear original. But it fails for equal families, where the makers declared no primary realization, and it imposes on those works a hierarchy they were designed to avoid.

The second takes the \emph{family} as the work: the set of all its realizations, each equally part of it. This respects equal families, but a set is not quite enough. Two unrelated films of equal length form a set. What makes a family a work is how its members relate.

The third takes the \emph{generative specification} as the work: the scene model, the rules for realizing it, and the transformations allowed. This fits works made by the virtual production of \Cref{ch:production}, where realizations are rendered from a model. It fits less well works shot on film, where the realizations are primary and no model exists apart from them.

The fourth takes the \emph{record of exhibition} as the work: every realization as actually shown, in every venue, at every screening. This is the most complete answer and the least usable. It makes the work unbounded and never finished, since each screening adds to it.

\section{Identity in constraints}

The book's answer combines the second and third. A selective work is identified by its \emph{declared constraints}: the shared world, the realizations authorized, the synchronization that ties them together, the transformations that relate them, and the rules governing how viewers may move among them.

The earlier chapters have already named these constraints. They are the spine and columns of the variant matrix in \Cref{ch:matrix}. They are the envelopes and rendezvous that keep the family synchronized. They are the switch system of \Cref{ch:switch-admissibility}, which states which histories can be continued into which realizations. They are the provenance of \Cref{ch:versions} and the registry of \Cref{ch:persuasion}. Gathered together, they can be written as
\begin{equation}
\mathfrak F=\big\langle\ \mathcal M,\ \mathcal V,\ \mathcal A,\ \mathcal S,\ \mathcal L,\ \mathcal E\ \big\rangle,
\end{equation}
where $\mathcal M$ is the shared world model, $\mathcal V$ the authorized realizations, $\mathcal A$ the synchronization anchors, $\mathcal S$ the selection and switching rules, $\mathcal L$ the transformations linking realizations, and $\mathcal E$ the exhibition conditions under which the work is meant to be shown.

A screening is a screening of the work when what it shows is generated within these constraints: realizations from the authorized set, synchronized at the declared anchors, switchable under the declared rules, exhibited under the declared conditions. A realization that violates them, an undeclared distributive alteration, an unauthorized stream, a sponsor's insertion not in the registry, is not the work, however it is titled.

\section{A precedent}

The idea that a work can be identified by its constraints rather than by a single instance has precedents. A chess problem is identified by its position and its stipulation, not by any particular game played from it; many solutions may be attempted, and the problem is the same problem. A piece of music with improvised passages is identified by its structure and its fixed parts, not by any single performance. A play is identified by its text, and every production realizes it differently.

Goodman gave the distinction a name. In his terms, an \emph{allographic} art is one whose works are identified by compliance with a notation, as a symphony is by its score, rather than by the history of a particular object, as a painting is \cite{goodman}. The variant matrix and switch system are not a notation in Goodman's strict technical sense, but they play a comparable role: they state what any authorized realization of the work must satisfy. Selective cinema moves film toward the allographic side of that line.

Cinema has been unusual in fixing its works to a single sequence of frames. Selective cinema returns it toward the condition of the other arts, in which a work is a structure that admits many realizations and remains itself across them. The Ship of Theseus puzzle arises because a ship seems to be its planks. A selective film is not its frames. It is the design that determines which frames may be shown together, to whom, and in what relation.

\section{What archives must keep}

The answer has practical consequences for preservation. An archive that keeps only one realization has kept a member of the work, not the work. An archive that keeps every realization but not their relations has kept a pile of films. What must be kept is the declared constraints and the realizations together: the variant matrix, the switch system, the registry and manifests, every authorized realization, the designed composite if there is one, the complementary score, and, for works produced from models, the models and the rules for rendering them, with the software needed to run them.

To these should be added representative records of exhibition: what was actually shown at some screenings, with their manifests. These are not the work, but they are evidence of how the work was realized, as recordings of performances are evidence of a piece of music.

\section{Authorship relocated}

The answer also says something about authorship. In ordinary cinema, the authors choose every frame. In selective cinema, they choose a space of frames and the rules for moving through it. Some of the authorship moves from the sequence to the design: from deciding what is seen to deciding what may be seen, in what combinations, by whom.

This does not diminish authorship. The designers of a variant family decide more, not less, than the makers of a single film: each realization, and every relation among them. What it changes is where authorship is visible. A single film displays its choices. A variant family displays some of them to each viewer and keeps the rest in its structure, which only the whole family, and the conversation of an audience that saw different parts of it, can reveal.
